
\begin{frame}{Gram-Schmidt}
  \begin{center}
    Em particular, não provaremos a equivalência entre os produtos internos, mas sim entre diferentes versões do algoritmo de Gram-Schmidt.
  \end{center}
\end{frame}

\begin{frame}[fragile]{Gram-Schmidt}
  \begin{minipage}[t]{0.45\textwidth}
    Gram-Schmidt Clássico:
    \vspace{0.2cm}
    \begin{code}
function Gram-Schmidt-C($A$)
    for $i$ in $1:n$
        $v_i$ = $a_i$
        for $j$ in $1:i-1$
            $v_i$ = $v_i - q_j q_j\t a_i$
        end
        $q_i$ = $v_i / \norm{v_i}$
    end
    return $Q$
end
    \end{code}
  \end{minipage}
  \pause
  \hfill
  \vrule
  \hfill
  \begin{minipage}[t]{0.50\textwidth}
    Gram-Schmidt Modificado:
    \vspace{0.2cm}
    \begin{code}
function Gram-Schmidt-M($A$)
    for $i$ in $1:n$
        $v_i$ = $a_i$
        for $j$ in $1:i-1$
            $v_i$ = $v_i - q_j q_j\t v_i$
        end
        $q_i$ = $v_i / \norm{v_i}$
    end
    return $Q$
end
    \end{code}
  \end{minipage}
  \vfill
\end{frame}

\begin{frame}[fragile]{Gram-Schmidt}
  \begin{minipage}[t]{0.45\textwidth}
    Gram-Schmidt Clássico:
    \vspace{0.2cm}
    \begin{code}
  for $j$ in $1:i-1$
      $v_i$ = $v_i - q_j q_j\t a_i$
  end
    \end{code}
  \end{minipage}
  \pause
  \hfill
  \vrule
  \hfill
  \begin{minipage}[t]{0.50\textwidth}
    Gram-Schmidt Modificado:
    \vspace{0.2cm}
    \begin{code}
  for $j$ in $1:i-1$
      $v_i$ = $v_i - q_j q_j\t v_i$
  end
    \end{code}
  \end{minipage}
  \vfill
  \pause
  \begin{minipage}[t]{0.45\textwidth}
    \[v_i = a_i - \sum_{j=1}^{i-1} q_j q_j\t a_i\]
  \end{minipage}
  \pause
  \hfill
  \vrule
  \hfill
  \begin{minipage}[t]{0.50\textwidth}
    \[v_i = \left(\prod_{j=i-1}^{1} \left(I - q_j q_j\t\right) \right) a_i\]
  \end{minipage}
  \vfill
\end{frame}

\begin{frame}{Gram-Schmidt}
  \begin{calculation}[\calciff]
    v_i = a_i - \sum_{j=1}^{i-1} q_j q_j\t a_i
    \step{Distributividade}
    v_i = \left(I - \sum_{j=1}^{i-1} q_j q_j\t \right) a_i
  \end{calculation}
\end{frame}

\begin{frame}{Gram-Schmidt}
  \begin{minipage}[t]{0.45\textwidth}
    \[v_i = \left(I - \sum_{j=1}^{i-1} q_j q_j\t \right) a_i\]
  \end{minipage}
  \pause
  \hfill
  \vrule
  \hfill
  \begin{minipage}[t]{0.50\textwidth}
    \[v_i = \left(\prod_{j=i-1}^{1} \left(I - q_j q_j\t\right) \right) a_i\]
  \end{minipage}
  \vfill
  \pause
  \begin{minipage}[t]{0.45\textwidth}
    \[I - \sum_{j=1}^{i-1} q_j q_j\t\]
  \end{minipage}
  \pause
  \hfill
  \vrule
  \hfill
  \begin{minipage}[t]{0.50\textwidth}
    \[\prod_{j=i-1}^{1} \left(I - q_j q_j\t\right)\]
  \end{minipage}
\end{frame}

\begin{frame}{Gram-Schmidt}
  \begin{center}
    É relativamente fácil mostrar por indução que essas duas matrizes são iguais sob a hipótese do algoritmo de Gram-Schmidt estar correto.
  \end{center}
  \pause
  \vspace{0.5cm}
  \begin{center}
    Mas qual é a graça? =)
  \end{center}
\end{frame}

\begin{frame}{Gram-Schmidt}
  Caso $i=1$:
  \begin{calculation}[=]
    I - \sum_{j=1}^{i-1} q_j q_j\t
    \step{$\displaystyle\sum_{j=1}^{0} q_j q_j\t = 0$}
    I
    \step{$\displaystyle\prod_{j=0}^{1} \left(I - q_j q_j\t\right) = I$}
    \prod_{j=i-1}^{1} \left(I - q_j q_j\t\right)
  \end{calculation}
\end{frame}

\begin{frame}{Gram-Schmidt}
  Hipótese de Indução:
  \begin{calculation}[=]
    I - \sum_{j=1}^{i-1} q_j q_j\t
    \step{Hipótese de Indução}
    \prod_{j=i-1}^{1} \left(I - q_j q_j\t\right)
  \end{calculation}
\end{frame}

\begin{frame}{Gram-Schmidt}
  Passo Indutivo:
  \begin{calculation}[=]
    \prod_{j=i}^{1} \left(I - q_j q_j\t\right)
    \step{Reindexação do Produtório}
    \left(I - q_{i} q_{i}\t\right) \prod_{j=i-1}^{1} \left(I - q_j q_j\t\right)
    \step{Hipótese de Indução}
    \left(I - q_{i} q_{i}\t\right) \left(I - \sum_{j=1}^{i-1} q_j q_j\t\right)
  \end{calculation}
\end{frame}

\begin{frame}{Gram-Schmidt}
  \vspace{-1cm}
  \begin{calculation}[=]
    \step{Distributividade}
    I^2 - \left(\sum_{j=1}^{i-1} q_j q_j\t\right) - q_{i} q_{i}\t + \left(\sum_{j=1}^{i-1} q_j q_j\t q_{i+1} q_{i+1}\t\right)
    \step{$q_j\t q_{i} = 0$ e $I^2 = I$}
    I - \left(\sum_{j=1}^{i-1} q_j q_j\t\right) - q_{i} q_{i}\t
    \step{Reindexação do Somatório}
    I - \sum_{j=1}^{i} q_j q_j\t
  \end{calculation}
  \qed
\end{frame}

\begin{frame}{Gram-Schmidt}
  \begin{center}
    Além de ser uma prova bem cuidadosa com os índices, perceba que precisamos assumir que $q_i\t q_j = 0$ para todo $i \neq j$.
  \end{center}
  \pause
  \vspace{0.3cm}
  \begin{center}
    Não só isso, mas, assim como no produto interno, tivemos que sair da implementação computacional e voltar para as definições algébricas dos algoritmos para estarmos em uma linguagem plausível para realizar a demonstração.
  \end{center}
  \pause
  \vspace{0.3cm}
  \begin{center}
    Reescreveremos esses algoritmos usando Stream Calculus!
  \end{center}
\end{frame}
